\documentclass[../main.tex]{subfiles}
\begin{document}

%========================================================================================
% FILENAME: 02_model.tex (v1.0.0)
%
% §2 Model — the objects the interface is built from: reserve asset and pool, the one
% conserved RECEIPT in two classes, the redemption-rate floor, the published constants and
% escrow, the two phases, genesis, and the state tuple.  No specified property proved here.
%
% Silences: no consumer noun (operator = the specification's admin role, introduced here); no
% realization reference.  Two clocks: cycle only.
%
% Change history lives in version control (see the versioning policy in main.tex).
%========================================================================================

\section{Model}
\label{sec:attestation:model}

This chapter fixes the objects the interface is built from: a reserve pool, one conserved
receipt in two classes, the redemption-rate floor that converts between them, the published
constants, and the two phases. No specified property is proved here; this chapter fixes the
state on which the rest of the document operates.

% ═══════════════════════════════════════════════════════════════════════
\subsection{Reserve Asset and Reserve Pool}
\label{sec:model:reserve}
% ═══════════════════════════════════════════════════════════════════════

\begin{definition}[Reserve Asset]
    \label{def:model:reserve-asset}
    A \emph{reserve asset} is a unit of value whose worth is independent of any system
    built on the receipt. All monetary quantities in this document are denominated in
    \emph{reserve units}.
\end{definition}

\begin{remark}[Instantiations]
    \label{rem:model:instantiations}
    The reserve asset is abstract; a deployment instantiates it (for example, L-BTC, or an
    external value token, EVT). The choice is a deployment matter, and nothing in this
    document depends on which instantiation is used.
\end{remark}

\begin{definition}[Reserve Pool]
    \label{def:model:reserve-pool}
    The \emph{reserve pool} $\reservepool{t} \in \R_{>0}$ is a publicly auditable, full
    reserve of the reserve asset held against the receipt. It rises on deposit and falls on
    redemption.
\end{definition}

% ═══════════════════════════════════════════════════════════════════════
\subsection{The Conserved Receipt}
\label{sec:model:claim}
% ═══════════════════════════════════════════════════════════════════════

\begin{definition}[Classes]
    \label{def:model:classes}
    The instrument issued under \emph{Attestation} is a single conserved
    \emph{reserve receipt} --- a redeemable claim
    on the reserve pool $\reservepool{}$ (\zcref{def:model:reserve-pool}), redeemable at the
    redemption-rate floor $\ratefloor{} = \reservepool{}/\totsupply{}$
    (\zcref{def:model:ratefloor}). The receipt is held in two \emph{redemption classes} ---
    classes, not assets:
    \begin{description}[nosep]
        \item[\textbf{Live class.}] Outstanding supply $\livesupply{}$. A \liveclass{} unit
            admits exactly three fates: \emph{traded} (transferred), \emph{burned}
            (\zcref{def:operations:burn}), or \emph{redeemed}
            (\zcref{def:operations:redemption}) for a floor of $\ratefloor{}$ reserve units
            per unit.
        \item[\textbf{Time-locked class.}] Outstanding supply $\tlocksupply{}$. A
            \tlockclass{} unit is transferable, but neither the burn predicate nor the
            redemption predicate is defined on it --- there is no operation consuming a
            \tlockclass{} unit into a burn record or a reserve payout --- so it is
            structurally non-burnable and non-redeemable, and carries no redemption floor.
            It exists only during the bootstrapping phase
            (\zcref{def:model:bootstrapping-phase}) and converts once into the live class at
            maturity (\zcref{postc:maturity:announcement,def:maturity:conversion}).
    \end{description}
    The single floor over the combined supply is the algebraic form of the two classes
    being one receipt: a distinct asset would carry its own floor
    (\zcref{def:model:ratefloor}).
\end{definition}

\begin{definition}[Supply]
    \label{def:model:supply}
    \[
        \totsupply{t} = \livesupply{t} + \tlocksupply{t}.
    \]
    The sum is conserved across the class boundary, so maturity conversion
    (\zcref{def:maturity:conversion}) leaves $\totsupply{}$ unchanged. Both terms count
    units of the one receipt.
\end{definition}

\begin{definition}[Redemption-Rate Floor]
    \label{def:model:ratefloor}
    \[
        \ratefloor{t} = \frac{\reservepool{t}}{\totsupply{t}},
    \]
    well-defined whenever $\totsupply{t} > 0$. There is one floor over the combined supply:
    \liveclass{} units are redeemed, and \tlockclass{} units valued, against the same pool
    $\reservepool{}$ at this single floor. The two classes therefore cannot be independent
    assets --- a distinct asset would carry its own floor; the shared $\ratefloor{}$ is the
    algebraic form of their being one receipt (\zcref{def:model:classes}). The floor only
    ever rises (\zcref{thm:invariants:floor-non-decrease}); it is a redemption floor, not a
    return.
\end{definition}

% ═══════════════════════════════════════════════════════════════════════
\subsection{Constants and Escrow}
\label{sec:model:constants}
% ═══════════════════════════════════════════════════════════════════════

\begin{definition}[Fee Rate and Split]
    \label{def:model:fee-and-split}
    The fee rate $\feerate \in (0,1)$ and the split parameter $\livesplit \in (0,1)$ are
    published constants of this layer, fixed at genesis. When units are minted, fraction
    $\livesplit$ enters the live class and $(1-\livesplit)$ the time-locked class during
    the bootstrapping phase (\zcref{def:model:bootstrapping-phase}); all into the live class
    during the normal phase (\zcref{def:model:normal-phase}). The \emph{operator} here is
    the specification-internal administrator role --- fee recipient and maturity announcer --- an internal
    concept of this layer.
    % operator privilege has a cross-system character priced downstream; not stated in this layer.
\end{definition}

\begin{remark}[Calibration Status]
    \label{rem:model:calibration}
    The numeric values carried for $\livesplit$ and $\feerate$ throughout this
    document ($\livesplit = \feerate = 0.5$) are \emph{placeholders}, not specified
    values. This layer fixes the two constants' roles and admissible ranges
    ($\livesplit, \feerate \in (0,1)$), not their magnitudes; the magnitudes are a
    calibration output, to be chosen \emph{jointly} by standard sensitivity-analysis
    methods. Jointly, because the exported quantities couple them: the net live share
    $\netliveshare$ (\zcref{def:interface:net-live-share}), the pre-deposit floor factor
    $\genesisfloorfactor = 1/(1-\livesplit)$
    (\zcref{cor:containment:pre-deposit-floor-factor}), the pre-deposit capacity
    $\attestcap = \genesisreserve\ln(1/(1-\livesplit))$
    (\zcref{prop:interface:bootstrap-capacity}), and the pre-maturity capacity bound
    $(\genesisreserve + D)\ln(1/(1-\livesplit))$
    (\zcref{thm:containment:pre-maturity-deposit-capacity}) are each functions of
    $\livesplit$ and $\feerate$ \emph{together}, so neither can be calibrated in
    isolation. ``Fixed at genesis'' (\zcref{def:model:fee-and-split})
    governs immutability \emph{after} the choice, not the choice itself; every numeric
    appearance below is illustrative (cf.\ \zcref{tab:verification:constants}).
\end{remark}

\begin{definition}[Deposit Escrow]
    \label{def:model:escrow}
    % [core] — specification-internal; never exported.
    $\escrow{t} \in \R_{\geq 0}$: reserve units in deposit escrow, awaiting the next cycle
    boundary (\zcref{rem:interface:clocks}).
\end{definition}

% ═══════════════════════════════════════════════════════════════════════
\subsection{Phases, Genesis, and State}
\label{sec:model:genesis}
% ═══════════════════════════════════════════════════════════════════════

\begin{definition}[Bootstrapping Phase]
    \label{def:model:bootstrapping-phase}
    Cycles $[0, \matcycle)$. Every deposit mints live-class and time-locked-class units in
    ratio $\livesplit : (1-\livesplit)$. The time-locked-class supply inflates the
    denominator $\totsupply{}$, bounding burn capacity and floor volatility.
\end{definition}

\begin{definition}[Normal Phase]
    \label{def:model:normal-phase}
    Cycles $[\matcycle, \infty)$. Every time-locked-class unit has converted into the live
    class (\zcref{def:maturity:conversion}); no time-locked-class unit is minted; the
    economy reduces to a single circulating class.
\end{definition}

\begin{postulate}[Genesis]
    \label{postc:model:genesis}
    \stackmarginnote{Note:}{The genesis split is what makes
    \zcref{prop:interface:bootstrap-capacity} a theorem: only $\livesplit\genesisreserve$
    live-class units exist to burn before external deposits.}%
    At initialization:
    \[
        \reservepool{0} = \genesisreserve > 0, \quad
        \livesupply{0} = \livesplit\genesisreserve, \quad
        \tlocksupply{0} = (1-\livesplit)\genesisreserve, \quad
        \totsupply{0} = \genesisreserve, \quad
        \ratefloor{0} = 1.
    \]
    The invariant $\totsupply{} > 0$ holds for all operational $t$.
\end{postulate}

\begin{remark}[State Tuple]
    \label{rem:model:state}
    The receipt state at cycle-time $t$ is
    \[
        \bigl(\reservepool{t},\ \livesupply{t},\ \tlocksupply{t},\
              \set{\attest{a}(t)}_{a \in \addrspace},\ \escrow{t}\bigr).
    \]
    The index $a$ is an opaque address (\zcref{def:interface:address-space}); the binding of
    addresses to any identity is the importing system's import-time interpretation. The
    attestation family is indexed by $\addrspace$ --- the burn record, not a stored
    attribute of any holder.
\end{remark}

%========================================================================================
% END OF FILE: 02_model.tex (v1.0.0)
%========================================================================================

\end{document}
